% v2.4.0 P2 figure batch 01: the first nine P2 UIDs in canonical sequence.
\documentclass[UTF8,a4paper,11pt,openany]{ctexbook}
\usepackage{../common/statlearnbook}
\SLSetBookMetadata{P2 绘图批次 01}{统计学习方法讲义 v2.4.0 绘图集成验证}{前四册九幅 P2 正式绘图源码局部编译}
\begin{document}
\pagestyle{empty}
\raggedbottom

\typeout{SLFIG-HARNESS-BEGIN FIG-P020-01}
% v2.6.0 figure UID: FIG-P020-01
% Image-informed reverse drawing: preserve the dependency graph while making
% the main chain, audit return, text hierarchy, and endpoint clearance explicit.
\tikzset{slfig-FIG-P020-01/.style={font=\fontsize{10.0pt}{12.0pt}\selectfont,every path/.append style={line cap=round,line join=round}}}
\pgfkeysifdefined{/tikz/alt}{}{\tikzset{alt/.code={}}}
\begin{figure}[H]\centering\label{schema:V1-C01:CH-17}
\begin{tikzpicture}[slfig-FIG-P020-01,
  node distance=6.2mm,
  stage/.style={rounded corners=1.7mm,draw=SLBlue,line width=.90pt,
    text width=29mm,minimum height=13.2mm,align=center,inner xsep=3.2pt,inner ysep=3pt,
    font=\fontsize{10.5pt}{12.6pt}\selectfont,text=SLInk},
  alt={对象—关系—结论：数学语言从对象声明到任务陈述的依赖关系。每一条箭头都表示右侧内容使用左侧定义。}
]
  \node[stage,fill=SLBlueBG] (object) {{\bfseries 对象声明}\\[-.6pt]{\fontsize{10.0pt}{12.0pt}\selectfont 集合、类型与维数}};
  \node[stage,fill=SLTealBG,right=of object] (relation) {{\bfseries 关系与映射}\\[-.6pt]{\fontsize{10.0pt}{12.0pt}\selectfont 定义域%
    \tikz[baseline=-.55ex]{
      \path[use as bounding box] (0,-.75mm) rectangle (7mm,.75mm);
      \draw[-{Stealth[length=1.55mm,width=1.05mm]},draw=SLInk,line width=.72pt]
        (1.05mm,0)--(5.95mm,0);
    }%
    值域}};
  \node[stage,fill=SLBlueBG,right=of relation] (logic) {{\bfseries 运算与逻辑}\\[-.6pt]{\fontsize{10.0pt}{12.0pt}\selectfont 复合、量词与约束}};
  \node[stage,fill=SLTealBG,right=of logic] (task) {{\bfseries 可核验任务}\\[-.6pt]{\fontsize{10.0pt}{12.0pt}\selectfont 输入、输出与判据}};
  \begin{scope}[on background layer]
    \node[fit=(object)(task),rounded corners=2.4mm,fill=SLBlue!3,draw=none,
      inner xsep=2.6mm,inner ysep=2.8mm] {};
  \end{scope}
  \foreach \a/\b in {object/relation,relation/logic,logic/task}
    \draw[-{Stealth[length=2.25mm,width=1.45mm]},draw=SLBlue,line width=.98pt,
      shorten <=1.5mm,shorten >=1.5mm]
      (\a.east)--(\b.west);
  \draw[-{Stealth[length=2.0mm,width=1.25mm]},draw=SLGray,line width=.58pt,
      dash pattern=on 2.4pt off 2.0pt,shorten <=1.5mm,shorten >=1.5mm]
    (task.south)--++(0,-9mm)-| node[pos=.50,below=2.0pt,fill=white,inner sep=1pt,
      font=\fontsize{10.0pt}{12.0pt}\selectfont,text=SLGray] {逆向核对：任务所用定义逐项返回检查}
    (object.south);
\end{tikzpicture}
\caption{数学语言从对象声明到任务陈述的依赖关系。每一条箭头都表示右侧内容使用左侧定义。}\label{fig:V1-C01-language-flow}
\end{figure}

\clearpage
\typeout{SLFIG-HARNESS-END FIG-P020-01}

\typeout{SLFIG-HARNESS-BEGIN FIG-P210-01}
% v2.3.1 figure UID: FIG-P210-01
% v2.6.0 Image-reference reverse redraw: reinforce split/tree hierarchy without changing geometry or topology.
\newcommand{\SLKdtreeTitleStyle}{\fontsize{10.5pt}{12.6pt}\selectfont\bfseries}
\tikzset{
  slfig-FIG-P210-01/.style={font=\fontsize{9.2pt}{11.0pt}\selectfont,every path/.append style={line cap=round,line join=round}},
  slfig-FIG-P210-01-axis/.style={-{Stealth[length=1.9mm,width=1.15mm]},draw=SLInk,line width=.50pt},
  slfig-FIG-P210-01-root-split/.style={draw=SLBlue,line width=1pt},
  slfig-FIG-P210-01-y-split/.style={draw=SLTeal,line width=.90pt,dash pattern=on 3.4pt off 2.1pt},
  slfig-FIG-P210-01-x-split/.style={draw=SLBlue,line width=.78pt},
  slfig-FIG-P210-01-axis-label/.style={font=\fontsize{9.2pt}{11pt}\selectfont},
  slfig-FIG-P210-01-split-label/.style={font=\fontsize{8.7pt}{10.4pt}\selectfont,fill=white,fill opacity=.94,text opacity=1,inner sep=.5pt},
  slfig-FIG-P210-01-point-label/.style={font=\fontsize{8.7pt}{10.4pt}\selectfont,fill=white,fill opacity=.94,text opacity=1,inner sep=.5pt},
  slfig-FIG-P210-01-note/.style={font=\fontsize{8.7pt}{10.4pt}\selectfont,text=SLGray}
}
\forestset{
  slfig-FIG-P210-01-tree/.style={
    for tree={draw=SLBlue,fill=SLBlueBG,rounded corners=1.5mm,line width=.78pt,
      inner xsep=4pt,inner ysep=3pt,
      edge={draw=SLBlue,line width=.78pt,-{Stealth[length=1.9mm,width=1.2mm]}}}
  }
}
\begin{figure}[htbp]
\centering
\begin{minipage}[c]{.47\linewidth}
\centering
{\SLKdtreeTitleStyle 样本空间与切分}\par\smallskip
\begin{tikzpicture}[slfig-FIG-P210-01,x=.46cm,y=.42cm]
  \draw[slfig-FIG-P210-01-axis] (0,0)--(10.5,0) node[right,slfig-FIG-P210-01-axis-label]{$x$};
  \draw[slfig-FIG-P210-01-axis] (0,0)--(0,8.2) node[above,slfig-FIG-P210-01-axis-label]{$y$};
  \draw[slfig-FIG-P210-01-root-split] (7,0)--(7,8)
    node[above,slfig-FIG-P210-01-split-label,text=SLBlue] {1：$x$};
  \draw[slfig-FIG-P210-01-y-split] (0,4)--(7,4)
    node[pos=.18,above,slfig-FIG-P210-01-split-label,text=SLTeal] {2：$y$};
  \draw[slfig-FIG-P210-01-y-split] (7,6)--(10,6)
    node[right,slfig-FIG-P210-01-split-label,text=SLTeal] {2：$y$};
  \draw[slfig-FIG-P210-01-x-split] (5,0)--(5,4)
    node[pos=.25,right,slfig-FIG-P210-01-split-label,text=SLBlue] {3：$x$};
  \draw[slfig-FIG-P210-01-x-split] (9,0)--(9,6)
    node[pos=.55,left,slfig-FIG-P210-01-split-label,text=SLBlue] {3：$x$};
  \foreach \x/\y/\lab in {2/3/A,4/7/B,5/4/C,7/2/D,8/1/E,9/6/F}{
    \fill[SLInk] (\x,\y) circle (2.2pt);
    \node[slfig-FIG-P210-01-point-label,anchor=south west] at (\x+.08,\y+.08) {$\lab(\x,\y)$};
  }
  \node[slfig-FIG-P210-01-note,anchor=west] at (.2,-1.0)
    {实线：$x$ 切分\quad 长虚线：$y$ 切分};
\end{tikzpicture}
\end{minipage}\hfill
\begin{minipage}[c]{.50\linewidth}
\centering
{\SLKdtreeTitleStyle 相同切分生成的 kd 树}\par\smallskip
\begin{forest}
  slfig tree,
  slfig-FIG-P210-01-tree,
  for tree={font=\footnotesize,l sep=8mm,s sep=4.5mm,minimum width=17mm},
  [{$D(7,2)$\\\textcolor{SLBlue}{\fontsize{9.2pt}{10.8pt}\selectfont 实线 $x$}}
    [{$C(5,4)$\\\textcolor{SLTeal}{\fontsize{9.2pt}{10.8pt}\selectfont 虚线 $y$}}
      [{$A(2,3)$}]
      [{$B(4,7)$}]
    ]
    [{$F(9,6)$\\\textcolor{SLTeal}{\fontsize{9.2pt}{10.8pt}\selectfont 虚线 $y$}}
      [{$E(8,1)$}]
    ]
  ]
\end{forest}
\par\vspace{1mm}
{\fontsize{8.7pt}{10.4pt}\selectfont 叶节点次序与左侧空间位置对应}
\end{minipage}
\caption{六点示例的一棵平衡kd树。节点同时保存样本与本层切分轴。}\label{fig:V2-C02-kd-tree}
\end{figure}

\clearpage
\typeout{SLFIG-HARNESS-END FIG-P210-01}

\typeout{SLFIG-HARNESS-BEGIN FIG-P324-01}
% v2.3.1 figure UID: FIG-P324-01
% Proposed post-v2.3.1 polish for FIG-P324-01: protect key-label size and stroke hierarchy.
\tikzset{slfig-FIG-P324-01/.style={font=\fontsize{9.2pt}{11.0pt}\selectfont,every path/.append style={line cap=round,line join=round}}}
\begin{figure}[H]
\centering
\begin{tikzpicture}[slfig-FIG-P324-01,
  dist/.style={slfig node,draw=SLBlue,fill=SLBlueBG,line width=.78pt,minimum width=25mm,text width=22mm},
  weak/.style={trapezium,trapezium left angle=75,trapezium right angle=105,
    draw=SLTeal,fill=SLTealBG,line width=.78pt,minimum width=25mm,minimum height=10mm,align=center},
  scalar/.style={circle,draw=SLGold,fill=SLGoldBG,line width=.82pt,minimum size=10mm,inner sep=1pt,align=center},
  ensemble/.style={slfig node,draw=SLBlue,fill=SLBlueBG,line width=.78pt,double,double distance=1.1pt,minimum width=30mm},
  err/.style={diamond,aspect=1.7,draw=SLGold,fill=SLGoldBG,line width=.82pt,align=center,inner sep=1.5mm},
]
  \node[font=\bfseries\fontsize{9.7pt}{11.4pt}\selectfont,text=SLBlue] at (0,2.75) {权重更新通道};
  \node[dist] (dm) at (-5.0,1.45) {样本分布\\$D_m$};
  \node[weak] (gm) at (-1.85,1.45) {训练\\$G_m$};
  \node[err] (em) at (1.15,1.45) {误差\\$e_m$};
  \node[dist] (dn) at (4.25,1.45) {更新分布\\$D_{m+1}$};
  \draw[slfig arrow,line width=.88pt] (dm)--(gm);
  \draw[slfig arrow accent,line width=.90pt] (gm)--(em);
  \draw[slfig arrow,line width=.88pt] (em)--(dn);

  \node[font=\bfseries\fontsize{9.7pt}{11.4pt}\selectfont,text=SLTeal] at (-4.25,-.35) {加法模型通道};
  \node[weak] (gcopy) at (-2.40,-1.35) {弱分类器\\$G_m$};
  \node[scalar] (alpha) at (.15,-1.35) {$\alpha_m$};
  \node[ensemble] (fm) at (3.30,-1.35) {集成模型\\$F_m=F_{m-1}+\alpha_mG_m$};
  \draw[slfig arrow accent,line width=.90pt] (gcopy) to[bend right=13] (fm);
  \draw[slfig arrow,line width=.84pt] (alpha)--(fm);
  \draw[slfig guide,line width=.52pt] (gm.south)--(gcopy.north);
  \draw[-{Stealth[length=1.9mm,width=1.2mm]},draw=SLGold,line width=.78pt]
    (em.south)--(alpha.north)
    node[midway,right=2.4mm,font=\fontsize{8.8pt}{10.2pt}\selectfont,text=SLGold,
      fill=none,inner sep=0pt] {$e_m\mapsto\alpha_m$};
  \draw[slfig return,draw=SLBlue!72,line width=.58pt]
    (dn.south) |- ([yshift=-7mm]dn.south) -| (gm.south);
  \node[anchor=north,font=\fontsize{8.5pt}{10.0pt}\selectfont,text=SLTextGray]
    at (2.85,-.15) {仅权重更新返回训练};
  \node[font=\fontsize{8.5pt}{10.0pt}\selectfont,text=SLTextGray] at (0,-2.35)
    {形状编码：分布 / 分类器 / 标量 / 集成};
\end{tikzpicture}
\caption{AdaBoost中样本权重$D_m$决定下一轮训练重点，分类器权重$\alpha_m$与弱分类器$G_m$直接进入加法模型；两类权重位于不同计算支路。}\label{fig:V3-C03-adaboost-loop}
\end{figure}

\clearpage
\typeout{SLFIG-HARNESS-END FIG-P324-01}

\typeout{SLFIG-HARNESS-BEGIN FIG-P412-01}
% v2.3.1 figure UID: FIG-P412-01
% Proposed post-v2.3.1 polish for FIG-P412-01: protect key-label size and stroke hierarchy.
\tikzset{slfig-FIG-P412-01/.style={font=\fontsize{9.2pt}{11.0pt}\selectfont,every path/.append style={line cap=round,line join=round}}}
\begin{figure}[H]
\centering
\begin{tikzpicture}[slfig-FIG-P412-01,
  dev/.style={slfig node,minimum width=27mm,text width=24mm,line width=.78pt,
    font=\fontsize{9.4pt}{11.2pt}\selectfont},
  final/.style={slfig node accent,minimum width=27mm,text width=24mm,line width=.82pt,
    font=\fontsize{9.4pt}{11.2pt}\selectfont},
  flow/.style={slfig arrow,line width=.90pt},
  feedback/.style={slfig return,draw=SLTextGray!72!SLRule,line width=.66pt},
]
  \node[dev] (task) at (-4.5,1.0) {任务定义};
  \node[dev] (family) at (-1.5,1.0) {候选模型族};
  \node[dev] (train) at (1.5,1.0) {训练候选};
  \node[dev] (valid) at (4.5,1.0) {验证选择};
  \draw[flow] (task)--(family);
  \draw[flow] (family)--(train);
  \draw[flow] (train)--(valid);
  \draw[feedback] (valid.north) |- ([yshift=8mm]valid.north) -| (family.north)
    node[pos=.48,above,font=\fontsize{9pt}{10.6pt}\selectfont,text=SLTextGray,
      fill=white,inner sep=1pt] {开发反馈：仅返回候选族};

  \node[diamond,aspect=1.8,draw=SLGold,fill=SLGoldBG,line width=.84pt,
        align=center,inner sep=1.3mm,font=\fontsize{9.4pt}{11.2pt}\selectfont]
        (freeze) at (4.5,-.45) {冻结};
  \node[final] (test) at (4.5,-2.0) {一次最终测试};
  \node[final] (report) at (1.5,-2.0) {报告结果};
  \draw[-{Stealth[length=2.25mm,width=1.4mm]},draw=SLGold,line width=.98pt] (valid)--(freeze);
  \draw[-{Stealth[length=2.25mm,width=1.4mm]},draw=SLTeal,line width=.98pt] (freeze)--(test);
  \draw[-{Stealth[length=2.25mm,width=1.4mm]},draw=SLTeal,line width=.98pt] (test)--(report);

  \node[draw=SLRed,fill=SLRedBG,rounded corners=1.6mm,line width=.78pt,
        minimum width=26mm,minimum height=10mm,align=center,
        font=\fontsize{9.2pt}{11pt}\selectfont] (locked) at (7.5,-2.0)
    {锁定测试集\\禁止开发期访问};
  \draw[-{Stealth[length=1.9mm,width=1.15mm]},draw=SLRed,line width=.72pt] (locked)--(test);
  \node[font=\fontsize{9pt}{10.6pt}\selectfont,text=SLRed,below=2mm of locked]
    {无任何返回候选族的路径};
\end{tikzpicture}
\caption{监督学习方法选择闭环。测试集不进入返回候选模型族的反馈回路。}\label{fig:V3-C07-selection-loop}
\end{figure}

\clearpage
\typeout{SLFIG-HARNESS-END FIG-P412-01}

\typeout{SLFIG-HARNESS-BEGIN FIG-P445-01}
% v2.3.1 figure UID: FIG-P445-01
% Proposed post-v2.3.1 polish for FIG-P445-01: protect key-label size and stroke hierarchy.
\tikzset{slfig-FIG-P445-01/.style={font=\fontsize{9.2pt}{11.0pt}\selectfont,every path/.append style={line cap=round,line join=round}}}
\begin{figure}[H]
\centering
\begin{tikzpicture}[slfig-FIG-P445-01,x=1.35cm,y=1.35cm,
  cluster band/.style={rounded corners=2.2pt,line width=.44pt,draw opacity=.30,fill opacity=.075},
  cluster label/.style={font=\bfseries\fontsize{9.6pt}{11.3pt}\selectfont},
  merge axis/.style={-{Stealth[length=2.0mm,width=1.18mm]},line width=.72pt},
  merge tick/.style={line width=.50pt},
  dendro branch/.style={draw=SLInk,line width=.98pt},
  cut line/.style={draw=SLGray!88!black,line width=.82pt,densely dashed},
]
  % Cluster bands lie behind branches and stop below the cut height 1.40.
  \path[cluster band,draw=SLBlue,fill=SLBlue] (-.35,.12) rectangle (1.35,1.34);
  \path[cluster band,draw=SLGold,fill=SLGold] (1.65,.12) rectangle (2.35,1.34);
  \path[cluster band,draw=SLTeal,fill=SLTeal] (2.65,.12) rectangle (4.35,1.34);
  \node[cluster label,text=SLBlue] at (.5,.25) {$C_1$};
  \node[cluster label,text=SLGold!88!black] at (2,.25) {$C_2$};
  \node[cluster label,text=SLTeal] at (3.5,.25) {$C_3$};

  \draw[merge axis] (-.65,.05)--(-.65,3.35)
    node[above,rotate=90,anchor=south,font=\fontsize{9.4pt}{11.2pt}\selectfont] {合并高度};
  \foreach \y in {0,1,2,3}{\draw[merge tick] (-.70,\y)--(-.60,\y) node[left=2pt,font=\fontsize{8.6pt}{10.2pt}\selectfont]{\y};}
  \foreach \x/\lab in {0/$x_1$,1/$x_2$,2/$x_3$,3/$x_4$,4/$x_5$}
    \node[font=\fontsize{9.4pt}{11.2pt}\selectfont,anchor=north] at (\x,0) {\lab};

  % Explicit merge heights: .65, .95, 2.10, 3.00.
  \draw[dendro branch] (0,.08)--(0,.65)--(1,.65)--(1,.08);
  \draw[dendro branch] (3,.08)--(3,.95)--(4,.95)--(4,.08);
  \draw[dendro branch] (3.5,.95)--(3.5,2.10)--(2,2.10)--(2,.08);
  \draw[dendro branch] (.5,.65)--(.5,3.00)--(2.75,3.00)--(2.75,2.10);
  \draw[cut line] (-.45,1.40)--(4.45,1.40);
  \node[font=\fontsize{9.2pt}{11pt}\selectfont,text=SLGray!88!black,anchor=west] at (4.50,1.40) {切割高度 $h_c=1.4$};
\end{tikzpicture}
\caption{树的纵坐标表示合并高度；在虚线高度切割时，虚线以下的连通分支分别成为一个类。}\label{fig:V4-C02-dendrogram}
\end{figure}

\clearpage
\typeout{SLFIG-HARNESS-END FIG-P445-01}

\typeout{SLFIG-HARNESS-BEGIN FIG-P467-01}
% v2.3.1 figure UID: FIG-P467-01
% Proposed post-v2.3.1 polish for FIG-P467-01: protect key-label size and stroke hierarchy.
\tikzset{slfig-FIG-P467-01/.style={font=\fontsize{9.2pt}{11.0pt}\selectfont,every path/.append style={line cap=round,line join=round}}}
\begin{figure}[H]
\centering
\begin{tikzpicture}[slfig-FIG-P467-01,
  curve trace/.style={draw=SLGray!80,line width=.52pt},
  principal vector/.style={-{Stealth[length=1.95mm,width=1.18mm]},draw=SLBlue,line width=.98pt},
  gray basis/.style={-{Stealth[length=1.72mm,width=1.05mm]},draw=SLGray!88!black,line width=.68pt},
  teal basis/.style={-{Stealth[length=1.72mm,width=1.05mm]},draw=SLTeal,line width=.68pt},
  phase arrow/.style={-{Stealth[length=2.1mm,width=1.25mm]},draw=SLGold,line width=.86pt},
]
\begin{groupplot}[
  group style={group size=4 by 1,horizontal sep=8mm},
  width=.20\linewidth,height=.20\linewidth,
  xmin=-1.9,xmax=1.9,ymin=-1.9,ymax=1.9,
  axis equal image,axis lines=middle,axis line style={draw=SLInk,line width=.58pt},
  xtick=\empty,ytick=\empty,clip=false,
  title style={font=\bfseries\fontsize{9.4pt}{11.2pt}\selectfont,yshift=1mm},
]
\nextgroupplot[title={输入：单位圆}]
  \addplot[curve trace,domain=0:360,samples=181] ({cos(x)},{sin(x)});
  \addplot[principal vector] coordinates {(0,0) (.8192,.5736)};
  \addplot[gray basis] coordinates {(0,0) (1,0)};
  \addplot[teal basis] coordinates {(0,0) (0,1)};
  \addplot[only marks,mark=*,mark size=1.35pt,draw=SLBlue!58,fill=SLBlue!58]
    coordinates {(.5,.866) (-.5,.866) (-.866,-.5) (.866,-.5)};
\nextgroupplot[title={$V^{\mathsf T}$：正交旋转}]
  \addplot[curve trace,domain=0:360,samples=181] ({cos(x)},{sin(x)});
  \addplot[principal vector] coordinates {(0,0) (.9962,.0872)};
  \addplot[gray basis] coordinates {(0,0) (.8660,-.5000)};
  \addplot[teal basis] coordinates {(0,0) (.5000,.8660)};
\nextgroupplot[title={$\Sigma$：轴向缩放}]
  \addplot[curve trace,domain=0:360,samples=181] ({1.6*cos(x)},{.65*sin(x)});
  \addplot[principal vector] coordinates {(0,0) (1.5939,.0567)};
  \addplot[gray basis] coordinates {(0,0) (1.3856,-.3250)};
  \addplot[teal basis] coordinates {(0,0) (.8000,.5629)};
\nextgroupplot[title={$U$：正交旋转}]
  \addplot[curve trace,domain=0:360,samples=181]
    ({1.4501*cos(x)-.2747*sin(x)},{.6762*cos(x)+.5891*sin(x)});
  \addplot[principal vector] coordinates {(0,0) (1.4205,.7248)};
  \addplot[gray basis] coordinates {(0,0) (1.3931,.2910)};
  \addplot[teal basis] coordinates {(0,0) (.4871,.8483)};
\end{groupplot}
\draw[phase arrow] (group c1r1.east)--(group c2r1.west);
\draw[phase arrow] (group c2r1.east)--(group c3r1.west);
\draw[phase arrow] (group c3r1.east)--(group c4r1.west);
\node[font=\fontsize{9pt}{10.6pt}\selectfont,text=SLGray!94!black] at ([yshift=-6mm]group c2r1.south east) {蓝向量与两条基方向贯穿四个面板；坐标范围一致};
\end{tikzpicture}
\caption{SVD把线性变换分成正交变换、轴向缩放和正交变换}\label{fig:V4-C03-svd-geometry}
\end{figure}

\clearpage
\typeout{SLFIG-HARNESS-END FIG-P467-01}

\typeout{SLFIG-HARNESS-BEGIN FIG-P504-01}
% v2.3.1 figure UID: FIG-P504-01
% Proposed post-v2.3.1 polish for FIG-P504-01: protect key-label size and stroke hierarchy.
\tikzset{
  slfig-FIG-P504-01/.style={font=\fontsize{9.4pt}{11.2pt}\selectfont,
    every path/.append style={line cap=round,line join=round}},
  slfig-FIG-P504-01-axis/.style={draw=SLGray,line width=.52pt,-{Stealth[length=1.6mm,width=1mm]}},
  slfig-FIG-P504-01-basis/.style={-{Stealth[length=2.1mm,width=1.25mm]},line width=1pt,draw=SLBlue},
  slfig-FIG-P504-01-target/.style={-{Stealth[length=2.1mm,width=1.25mm]},line width=1.02pt,draw=SLGold},
  slfig-FIG-P504-01-residual/.style={draw=SLGray,densely dashed,line width=.58pt},
  slfig-FIG-P504-01-right-angle/.style={draw=SLBlue,line width=.68pt},
  slfig-FIG-P504-01-reconstruction/.style={-{Stealth[length=1.9mm,width=1.15mm]},draw=SLTeal,line width=.96pt},
  slfig-FIG-P504-01-title/.style={font=\bfseries\fontsize{10.2pt}{12.2pt}\selectfont},
  slfig-FIG-P504-01-note/.style={font=\fontsize{9.2pt}{11pt}\selectfont,align=center,text=SLGray},
  slfig-FIG-P504-01-summary/.style={font=\fontsize{9.2pt}{11pt}\selectfont,
    draw=SLRuleGray,line width=.58pt,rounded corners=2pt,fill=SLSoftGray,
    minimum width=116mm,minimum height=6.5mm,align=center},
}
\begin{figure}[H]
\centering
\begin{tikzpicture}[slfig-FIG-P504-01,>=Stealth,
  every node/.style={font=\fontsize{9.4pt}{11.2pt}\selectfont},
]
  % 两面板共享同一目标向量 x 与秩 K=2。
  \begin{scope}[shift={(-3.55,0)}]
    \node[slfig-FIG-P504-01-title,text=SLBlue] at (0,2.62) {LSA：正交子空间（$K=2$）};
    \fill[SLBlue!6] (-2.15,-1.55) rectangle (2.15,1.70);
    \draw[slfig-FIG-P504-01-axis] (-2.35,0)--(2.45,0);
    \draw[slfig-FIG-P504-01-axis] (0,-1.75)--(0,1.95);
    \draw[slfig-FIG-P504-01-basis] (0,0)--(1.65,0) node[below right] {$u_1$};
    \draw[slfig-FIG-P504-01-basis] (0,0)--(0,1.45) node[above left] {$u_2$};
    \draw[slfig-FIG-P504-01-right-angle] (.23,0)--(.23,.23)--(0,.23);
    \draw[slfig-FIG-P504-01-target] (0,0)--(-1.05,1.28) node[above left] {$x$};
    \draw[slfig-FIG-P504-01-residual] (-1.05,1.28)--(-1.05,0);
    \fill[SLGold] (-1.05,0) circle (1.85pt);
    \node[anchor=north,text=SLGold] at (-1.05,-.06) {$\widehat x=U_2U_2^{\mathsf T}x$};
    \node[slfig-FIG-P504-01-note] at (0,-2.05)
      {$y=U_2^{\mathsf T}x$，坐标可正可负};
  \end{scope}

  \begin{scope}[shift={(3.55,0)}]
    \node[slfig-FIG-P504-01-title,text=SLTeal] at (0,2.62) {NMF：非负锥（$K=2$）};
    \fill[SLTeal,opacity=.09] (0,0)--(2.10,.48)--(.72,1.95)--cycle;
    \draw[slfig-FIG-P504-01-axis] (-2.35,0)--(2.45,0);
    \draw[slfig-FIG-P504-01-axis] (0,-1.75)--(0,1.95);
    \draw[slfig-FIG-P504-01-basis,draw=SLTeal] (0,0)--(2.10,.48) node[right] {$w_1$};
    \draw[slfig-FIG-P504-01-basis,draw=SLTeal] (0,0)--(.72,1.95) node[above] {$w_2$};
    % 与左图完全相同的目标向量坐标。
    \draw[slfig-FIG-P504-01-target] (0,0)--(-1.05,1.28) node[above left] {$x$};
    \draw[slfig-FIG-P504-01-residual] (-1.05,1.28)--(.72,1.20);
    \draw[slfig-FIG-P504-01-reconstruction]
      (0,0)--(.72,1.20) node[right] {$Wh$};
    \fill[SLTeal] (.72,1.20) circle (1.85pt);
    \node[slfig-FIG-P504-01-note] at (0,-2.05)
      {$h_1,h_2\ge0$，只允许锥内组合};
  \end{scope}

  \node[slfig-FIG-P504-01-summary] at (0,-2.82)
    {同秩 $K=2$\quad/\quad 不同约束与目标：正交投影 vs. 非负重构};
\end{tikzpicture}
\caption{LSA的正交子空间表示与NMF的非负锥表示：二者维数相同，坐标约束和最优性不同}\label{fig:V4-C05-two-geometries}
\end{figure}

\clearpage
\typeout{SLFIG-HARNESS-END FIG-P504-01}

\typeout{SLFIG-HARNESS-BEGIN FIG-P521-01}
% v2.3.1 figure UID: FIG-P521-01
% v2.6.0 reverse redraw: clarify the generating chain, node classes, and nested-plate labels.
\tikzset{slfig-FIG-P521-01/.style={font=\fontsize{9.4pt}{11.2pt}\selectfont,every path/.append style={line cap=round,line join=round}}}
\begin{figure}[H]
\centering
\begin{tikzpicture}[slfig-FIG-P521-01,>=Stealth,
  every node/.style={font=\fontsize{9.4pt}{11.2pt}\selectfont,inner sep=1.8pt},
  rv/.style={circle,draw=SLBlue,line width=1.0pt,minimum size=9.8mm,fill=white},
  observed/.style={rv,fill=SLBlue!15},
  param/.style={rectangle,draw=SLTeal,line width=.82pt,rounded corners=2pt,
    fill=SLTeal!7,minimum width=18mm,minimum height=7.4mm,
    inner xsep=2.3mm,inner ysep=1.0mm},
  gen/.style={-{Stealth[length=2.2mm]},line width=1.0pt,draw=SLBlue},
  paramedge/.style={-{Stealth[length=2.0mm]},line width=.88pt,draw=SLTeal},
  plate/.style={draw=SLGray,line width=.62pt,rounded corners=2.3pt},
  platelabel/.style={text=SLTextGray,font=\fontsize{8.8pt}{10.4pt}\selectfont},
  legend/.style={anchor=west,text=SLTextGray,font=\fontsize{8.8pt}{10.4pt}\selectfont},
  summary/.style={draw=SLRuleGray,line width=.62pt,rounded corners=2.3pt,
    fill=SLSoftGray,align=center,minimum width=42mm,minimum height=11mm},
]
  \node[param] (theta) at (-2.85,1.75) {$\theta_d=P(z\mid d)$};
  \node[param] (phi) at (2.70,1.75) {$\phi_z=P(w\mid z)$};
  \node[observed] (d) at (-3.45,0) {$d$};
  \node[rv] (z) at (0,0) {$z$};
  \node[observed] (w) at (3.15,0) {$w$};
  \draw[gen] (d)--(z);
  \draw[gen] (z)--(w);
  \draw[paramedge] (theta)--(z);
  \draw[paramedge] (phi)--(w);

  % 内层词位 plate 与外层文档 plate。
  \draw[plate] (-.75,-.75) rectangle (3.88,.76);
  \node[platelabel,anchor=north] at (2.00,-.80) {$n=1{:}N_d$};
  \draw[plate] (-4.18,-1.25) rectangle (4.38,2.38);
  \node[platelabel,anchor=south east] at (4.33,-1.20) {$d=1{:}D$};

  \node[summary] at (0,-2.02)
    {$d\longrightarrow z\longrightarrow w$\qquad $w\perp d\mid z$};
  \node[legend] at (4.72,.55)
    {实心：观测};
  \node[legend] at (4.72,.05)
    {空心：潜变量};
  \node[legend] at (4.72,-.45)
    {矩形：参数};
\end{tikzpicture}
\caption{PLSA生成图：文档决定主题混合，主题决定单词分布；给定主题后单词与文档条件独立}\label{fig:V4-C06-plsa-dag}
\end{figure}

\clearpage
\typeout{SLFIG-HARNESS-END FIG-P521-01}

\typeout{SLFIG-HARNESS-BEGIN FIG-P525-01}
% v2.6.0 figure UID: FIG-P525-01
\tikzset{
  slfig-FIG-P525-01/.style={font=\fontsize{9.4pt}{11.2pt}\selectfont,
    every path/.append style={line cap=round,line join=round}},
  slfig-FIG-P525-01-outer/.style={draw=SLTextGray,line width=.92pt},
  slfig-FIG-P525-01-hull/.style={draw=SLTeal,line width=.86pt},
  slfig-FIG-P525-01-spoke/.style={draw=SLGray,densely dashed,line width=.54pt},
  slfig-FIG-P525-01-weight/.style={fill=none,inner sep=0pt},
  slfig-FIG-P525-01-doclabel/.style={anchor=west,text=SLGold,fill=white,
    fill opacity=.92,text opacity=1,inner xsep=1pt,inner ysep=.5pt},
  slfig-FIG-P525-01-formula/.style={draw=SLRuleGray,rounded corners=2pt,
    fill=SLSoftGray,line width=.62pt,align=center,inner ysep=2.2pt},
  slfig-FIG-P525-01-legend/.style={anchor=west,text=SLTextGray,
    font=\fontsize{8.8pt}{10.4pt}\selectfont}
}
\begin{figure}[H]
\centering
\begin{tikzpicture}[slfig-FIG-P525-01,x=1cm,y=1cm]
  % 显式重心变换：(p_1,p_2,p_3) ->
  % (x,y)=(6(p_2+p_3/2), 6(sqrt(3)/2)p_3)。
  \newcommand{\barypoint}[4]{%
    \coordinate (#1) at ({6*(#3+0.5*#4)},{6*0.8660254*#4});}
  \coordinate (Wone) at (0,0);
  \coordinate (Wtwo) at (6,0);
  \coordinate (Wthree) at (3,{3*sqrt(3)});
  \barypoint{T1}{.78}{.12}{.10}
  \barypoint{T2}{.10}{.78}{.12}
  \barypoint{T3}{.12}{.13}{.75}
  \barypoint{Doc}{.309}{.4195}{.2715}

  \fill[SLTeal,opacity=.085] (T1)--(T2)--(T3)--cycle;
  \draw[slfig-FIG-P525-01-outer] (Wone)--(Wtwo)--(Wthree)--cycle;
  \node[below left] at (Wone) {$w_1$};
  \node[below right] at (Wtwo) {$w_2$};
  \node[above] at (Wthree) {$w_3$};
  \draw[slfig-FIG-P525-01-hull] (T1)--(T2)--(T3)--cycle;

  \foreach \p/\lab/\pos in {T1/$\phi_{:1}$/left,T2/$\phi_{:2}$/right,T3/$\phi_{:3}$/above right}{
    \fill[SLBlue] (\p) circle (2.6pt) node[\pos,text=SLBlue] {\lab};
  }
  \draw[slfig-FIG-P525-01-spoke] (Doc)--(T1);
  \draw[slfig-FIG-P525-01-spoke] (Doc)--(T2);
  \draw[slfig-FIG-P525-01-spoke] (Doc)--(T3);
  \node[slfig-FIG-P525-01-weight] at (2.05,1.28) {$\theta_{1j}$};
  \node[slfig-FIG-P525-01-weight] at (3.65,.95) {$\theta_{2j}$};
  \node[slfig-FIG-P525-01-weight] at (2.55,2.05) {$\theta_{3j}$};
  \node[diamond,draw=SLGold,fill=SLGold!18,minimum size=6mm,
    inner sep=0pt,line width=.9pt] (dj) at (Doc) {};
  \node[slfig-FIG-P525-01-formula,
    minimum width=58mm,minimum height=8mm] at (8.0,2.55)
    {$P(w\mid d_j)=\sum_{k=1}^{3}\theta_{kj}\phi_{:k}$\\
     $\theta_{kj}\ge0,\ \sum_k\theta_{kj}=1$};
  \node[slfig-FIG-P525-01-legend] at (6.95,1.25)
    {圆：主题点\quad 菱形：文档分布 $P(w\mid d_j)$};
\end{tikzpicture}
\caption{三个单词维度中的主题单纯形：文档分布是主题点按$\theta_{:j}$形成的凸组合}\label{fig:V4-C06-simplex}
\end{figure}

\clearpage
\typeout{SLFIG-HARNESS-END FIG-P525-01}

\end{document}
